\documentclass[12pt, twoside]{article}
%\documentclass[12pt, twoside]{article}
\usepackage[letterpaper, margin=1in, headsep=0.2in]{geometry}
\setlength{\headheight}{0.6in}
%\usepackage[english]{babel}
\usepackage[utf8]{inputenc}
\usepackage{microtype}
\usepackage{amsmath}
\usepackage{amssymb}
%\usepackage{amsfonts}
\usepackage[nomessages]{fp} %\FPeval{\var-name}{2*sin(pi/6)}
\usepackage{siunitx} %units in math. eg 20\milli\meter
\usepackage{yhmath} % for arcs, overparenth command
\usepackage{tikz} %graphics
\usetikzlibrary{quotes, angles, arrows, arrows.meta}
\usepackage{pgfplots}
\usepgfplotslibrary{fillbetween}
\pgfplotsset{compat=1.16}
\usepackage{graphicx} %consider setting \graphicspath{{images/}}
\usepackage{parskip} %no paragraph indent
\usepackage{enumitem}
\usepackage{multicol}
\usepackage{venndiagram}

\usepackage{fancyhdr}
\pagestyle{fancy}
\fancyhf{}
\renewcommand{\headrulewidth}{0pt} % disable the underline of the header
\raggedbottom
\hfuzz=2mm %suppresses overfull box warnings

\usepackage{hyperref}

\fancyhead[RO]{\fbox{\begin{minipage}{6cm}\footnotesize Nickname:\\[0.5cm]\end{minipage}}}
\fancyhead[LE]{\fbox{\begin{minipage}{6cm}\footnotesize Nickname:\\[0.5cm]\end{minipage}}}
\fancyfoot[C]{\thepage}

\begin{document}
\fancyhead[LO]{La Scuola d'Italia / Dr.\ Huson / IB Math 11 (a) \\* 28 October 2026}
\subsubsection*{2.6 Early Finishers: Lines, Exponents, Radicals and Sequences}

\emph{For students who finish the quiz early. A calculator is allowed.}

\begin{enumerate}[itemsep=0.15cm]

%--- Problem 1: Same-base equation; integers a, b from a surd product; equation quadratic in 3^x [11 marks] ---
%--- Source: copy of 2-6Q_Lines-exponents-radicals version (b) P3. Original (AA SL 1.5, 1.7, 2.10); harder parallel of 2-6Q version (a) P3-P4 ---
\item[1.] [Maximum mark: 11]

\begin{enumerate}[itemsep=0.1cm]
  \item Solve the equation $25^{x + 1} = 125^{2x - 3}$. \hfill [3]
  \item Find the integers $a$ and $b$ such that
  $\left(2 + \sqrt{3}\right)\left(a + b\sqrt{3}\right) = 12 + 7\sqrt{3}$.
  \hfill [3]
\end{enumerate}

Consider the equation $3\left(9^{x}\right) - 82\left(3^{x}\right) + 27 = 0$.

\begin{enumerate}[itemsep=0.1cm]
  \item[(c)] By letting $u = 3^{x}$, solve the equation. \hfill [5]
\end{enumerate}

\begin{tikzpicture}
  \draw (-0.6,0) rectangle (14.6,15.9);
  \foreach \y in {1,2,...,15} \draw[dotted] (0,\y)--(14.1,\y);
\end{tikzpicture}

\newpage

%--- Problem 2: Arithmetic vs geometric: expressions, least n with v_n > u_n, least n with geometric sum > arithmetic sum [8 marks] ---
%--- Source: copy of 2-6Q_Lines-exponents-radicals version (b) P4. Original (AA SL 1.2, 1.3); unit 1 material, harder than 1-10Q_Sequences P4 ---
\item[2.] [Maximum mark: 8]

An arithmetic sequence has first term $u_{1} = 7$ and common difference
$3$. A geometric sequence has first term $v_{1} = 2$ and common ratio
$1.5$.

\begin{enumerate}[itemsep=0.1cm]
  \item Write down an expression for $u_{n}$ and an expression for $v_{n}$.
  \hfill [2]
  \item Find the least value of $n$ for which $v_{n} > u_{n}$. \hfill [3]
  \item Find the least value of $n$ for which the sum of the first $n$ terms
  of the geometric sequence is greater than the sum of the first $n$ terms
  of the arithmetic sequence. \hfill [3]
\end{enumerate}

\begin{tikzpicture}
  \draw (-0.6,0) rectangle (14.6,15.9);
  \foreach \y in {1,2,...,15} \draw[dotted] (0,\y)--(14.1,\y);
\end{tikzpicture}

\end{enumerate}

\newpage

\fancyhead[LO]{La Scuola d'Italia / Dr.\ Huson / IB Math 11 (b) \\* 28 October 2026}
\subsubsection*{2.6 Early Finishers: Lines, Exponents, Radicals and Sequences}

\emph{For students who finish the quiz early. A calculator is allowed.}

\begin{enumerate}[itemsep=0.15cm]

%--- Problem 1: Lines with a parameter k: parallel (two values), perpendicular, intersection when one line is vertical [10 marks] ---
%--- Source: Original (AA SL 2.1, 2.7); early-finisher extension, version (b) ---
\item[1.] [Maximum mark: 10]

The lines $L_{1}$ and $L_{2}$ have equations
$$kx + 2y = 7 \qquad \text{and} \qquad 3x + (k - 1)y = 5,$$
where $k$ is a real constant.

\begin{enumerate}[itemsep=0.1cm]
  \item Find the two values of $k$ for which $L_{1}$ and $L_{2}$ are
  parallel. \hfill [4]
  \item Find the value of $k$ for which $L_{1}$ and $L_{2}$ are
  perpendicular. \hfill [3]
  \item Let $k = 1$. Find the coordinates of the point where $L_{1}$ and
  $L_{2}$ intersect. \hfill [3]
\end{enumerate}

\begin{tikzpicture}
  \draw (-0.6,0) rectangle (14.6,14.0);
  \foreach \y in {1,2,...,13} \draw[dotted] (0,\y)--(14.1,\y);
\end{tikzpicture}

\newpage

%--- Problem 2: Geometric sequence with terms in k: quadratic for k with surd roots, r in surd form, sum to infinity [12 marks] ---
%--- Source: Original (AA SL 1.3, 1.8, 1.7); early-finisher extension, version (b); unit 1 material with surds ---
\item[2.] [Maximum mark: 12]

The first three terms of a geometric sequence are $k - 3$, $k + 1$ and
$3k - 1$.

\begin{enumerate}[itemsep=0.1cm]
  \item Show that $k^{2} - 6k + 1 = 0$. \hfill [2]
  \item Find the two possible values of $k$. Give your answers in the form
  $a + b\sqrt{2}$, where $a, b \in \mathbb{Z}$. \hfill [2]
  \item For the larger value of $k$, find the common ratio. Give your
  answer in the form $a + b\sqrt{2}$. \hfill [4]
  \item For the smaller value of $k$, show that the series has a sum to
  infinity, and find its value. \hfill [4]
\end{enumerate}

\begin{tikzpicture}
  \draw (-0.6,0) rectangle (14.6,15.0);
  \foreach \y in {1,2,...,14} \draw[dotted] (0,\y)--(14.1,\y);
\end{tikzpicture}

\end{enumerate}
\end{document}
